\section{Methodology}

\subsection{Problem Formulation}
We formulate the task as abstractive question answering: given a medical question $q$ and a relevant context passage $c$, the model generates an answer $a$ in natural Vietnamese. Formally, we train a sequence-to-sequence model $f_\theta$ such that:
\begin{equation}
  a = f_\theta(\texttt{"question: } q \texttt{ context: } c\texttt{"})
\end{equation}

\subsection{Model Architectures}
We fine-tune two Vietnamese Transformer-based Seq2Seq models:
\begin{itemize}
  \item \textbf{ViT5-base} \cite{phan-etal-2022-vit5}: A T5 architecture pretrained on Vietnamese text ($\sim$270M parameters).
  \item \textbf{BARTpho-word} \cite{tran-etal-2021-bartpho}: A BART architecture with word-level tokenization ($\sim$396M parameters).
\end{itemize}
A zero-shot \textbf{Llama-3.3-70B-versatile} baseline via Groq API is included as an external reference point.

\subsection{Retrieval-Augmented Generation (RAG)}
\label{sec:rag}
To improve the practical usability of our system, we integrate a \textbf{Retrieval-Augmented Generation (RAG)} pipeline at inference time. This allows end users to input only a question, without needing to provide a medical context document manually.

The RAG pipeline operates as follows:
\begin{enumerate}
  \item \textbf{Knowledge Base Construction:} We extract all unique context passages from the ViMedAQA dataset to form a medical corpus.
  \item \textbf{BM25 Retrieval:} Given a user question, we use BM25Okapi (a TF-IDF-based ranking function) to retrieve the top-1 most relevant context from the corpus.
  \item \textbf{Answer Generation:} The retrieved context is concatenated with the question in the training format (\texttt{question: \{q\} context: \{c\}}) and passed to the fine-tuned ViT5 model.
\end{enumerate}

\noindent\textbf{Note:} During training, models are fine-tuned with gold (ground-truth) contexts from the dataset. The BM25 retrieval module is used only at inference time in the web application, ensuring that model evaluation remains fair and comparable.

\subsection{Evaluation Metrics}
We evaluate generated answers using:
\begin{itemize}
  \item \textbf{ROUGE-1/2/L} — measures n-gram overlap between generated and reference answers.
  \item \textbf{BLEU-4} — measures 4-gram precision with brevity penalty.
  \item \textbf{BERTScore-F1} — measures semantic similarity using contextualized embeddings.
\end{itemize}
